% Generated draft: requires mathematical and source review.
\documentclass[11pt]{article}
\usepackage[margin=25mm]{geometry}
\usepackage{fontspec}
\setmainfont{Latin Modern Roman}
\usepackage{amsmath,amssymb,mathtools,unicode-math}
\setmathfont{Latin Modern Math}
\usepackage{enumitem,needspace,xurl,hyperref}
\setlength{\parindent}{0pt}
\setlength{\parskip}{5pt}
\setlength{\emergencystretch}{3em}
\newcommand{\N}{\mathbb{N}}
\newcommand{\Z}{\mathbb{Z}}
\newcommand{\Q}{\mathbb{Q}}
\newcommand{\R}{\mathbb{R}}
\newcommand{\C}{\mathbb{C}}
\newcommand{\F}{\mathbb{F}}
\newcommand{\A}{\mathbb{A}}
\newcommand{\PP}{\mathbb{P}}
\newcommand{\E}{\mathbb{E}}
\newcommand{\eps}{\varepsilon}
\newcommand{\dd}{\,\mathrm{d}}
\begin{document}
\section*{UnsolvedMath ID: 30000954}
\textbf{Source problem code:} catalogue entry
\subsection{Definitions and mathematical statement}
The source concerns a vector-valued dislocation phase-field energy, denoted here by $E_\varepsilon^{\Gamma}$, whose expected limit is a line-tension functional. The symbol $\Gamma(z)$ denotes an interaction kernel; it is distinct from the use of $\Gamma$ in $\Gamma$-convergence. The flat-interface upper bound is written in the source as
\[
F_{\mathrm{flat}}(u)=\int\gamma_0(\nu,[u])\,d\mathcal H^1.
\]
In the usual jump-set interpretation, this means integration on $J_u$, with $\nu=\nu_u$ a jump normal and $[u]=u^+-u^-$ the jump of a vector-valued $BV$ function. The domain, component number, allowed values and class of flat-interface competitors are not given in the record, so this interpretation does not specify those missing model data.

For a specified convergence topology $\tau$ on the intended $BV$ domain and the specified admissible class, the sequential lower-semicontinuous relaxation is
\[
\overline F_{\mathrm{flat}}^{\,\tau}(u)
=\inf\left\{\liminf_{j\to\infty}F_{\mathrm{flat}}(u_j):u_j\xrightarrow{\tau}u\right\},
\]
where $F_{\mathrm{flat}}$ is extended by $+\infty$ outside its admissible domain. The requested identification is
\[
\Gamma\text{-}\lim_{\varepsilon\downarrow0}E_\varepsilon^{\Gamma}
=\overline F_{\mathrm{flat}}^{\,\tau}.
\]
This would require the corresponding lower bound for every converging sequence and recovery sequences at every admissible limit, in the chosen topology and energy normalization.

The source-verification note says that the report already proves this identification for kernels
\[
\Gamma(z)=|z|^{-3}\gamma(z/|z|)
\]
with bounded uniformly elliptic angular part. The remaining problem is extension beyond that structural class, not reproving its stated restricted theorem. For a matrix angular part, uniform ellipticity conventionally entails fixed positive upper and lower quadratic-form bounds, but its exact symmetry, dimensions and hypotheses are not supplied. Nor is the larger admissible kernel class specified. The request must not be converted into an assertion for every arbitrary kernel.

The actual phase-field energy, its scaling and normalization, the formula for $\gamma_0$, boundary conditions, admissible $BV$ functions and topology are absent. In particular, ``in $BV$'' does not by itself choose between $L^1$, weak-star $BV$ or stronger convergence. The supplied assessment dated 2026-08-21 labels the record partially solved and emphasizes the restricted proven kernel result.
\subsection{Short English statement}
Beyond the report's homogeneous uniformly elliptic kernels, ask whether the phase-field $\Gamma$-limit equals the relaxation of the flat-interface line energy. The restricted result is reported proved; the broader kernel class, energy and relaxation topology remain unspecified.
\subsection{Sources}
\begin{itemize}
\item Local UnsolvedMath registry, permanent ID 30000954, source code catalogue entry; supplied complete source record. Its corrected-verified note explicitly identifies the open part as extension beyond a kernel class already treated in the report.
\item Atomistic Models of Materials: Mathematical Challenges (2009), Oberwolfach Reports, DOI 10.4171/owr/2008/21, pp. 1079--1156: \url{https://doi.org/10.4171/owr/2008/21}. Proposed-year metadata gives 2008.
\end{itemize}
\end{document}
